% Options for packages loaded elsewhere
% Options for packages loaded elsewhere
\PassOptionsToPackage{unicode}{hyperref}
\PassOptionsToPackage{hyphens}{url}
\PassOptionsToPackage{dvipsnames,svgnames,x11names}{xcolor}
%
\documentclass[
  nswissgerman,
  11pt,
  a4paper,
  DIV=11,
  numbers=noendperiod]{scrartcl}
\usepackage{xcolor}
\usepackage[margin=2.5cm]{geometry}
\usepackage{amsmath,amssymb}
\setcounter{secnumdepth}{5}
\usepackage{iftex}
\ifPDFTeX
  \usepackage[T1]{fontenc}
  \usepackage[utf8]{inputenc}
  \usepackage{textcomp} % provide euro and other symbols
\else % if luatex or xetex
  \usepackage{unicode-math} % this also loads fontspec
  \defaultfontfeatures{Scale=MatchLowercase}
  \defaultfontfeatures[\rmfamily]{Ligatures=TeX,Scale=1}
\fi
\usepackage{lmodern}
\ifPDFTeX\else
  % xetex/luatex font selection
\fi
% Use upquote if available, for straight quotes in verbatim environments
\IfFileExists{upquote.sty}{\usepackage{upquote}}{}
\IfFileExists{microtype.sty}{% use microtype if available
  \usepackage[]{microtype}
  \UseMicrotypeSet[protrusion]{basicmath} % disable protrusion for tt fonts
}{}
\usepackage{setspace}
\makeatletter
\@ifundefined{KOMAClassName}{% if non-KOMA class
  \IfFileExists{parskip.sty}{%
    \usepackage{parskip}
  }{% else
    \setlength{\parindent}{0pt}
    \setlength{\parskip}{6pt plus 2pt minus 1pt}}
}{% if KOMA class
  \KOMAoptions{parskip=half}}
\makeatother
% Make \paragraph and \subparagraph free-standing
\makeatletter
\ifx\paragraph\undefined\else
  \let\oldparagraph\paragraph
  \renewcommand{\paragraph}{
    \@ifstar
      \xxxParagraphStar
      \xxxParagraphNoStar
  }
  \newcommand{\xxxParagraphStar}[1]{\oldparagraph*{#1}\mbox{}}
  \newcommand{\xxxParagraphNoStar}[1]{\oldparagraph{#1}\mbox{}}
\fi
\ifx\subparagraph\undefined\else
  \let\oldsubparagraph\subparagraph
  \renewcommand{\subparagraph}{
    \@ifstar
      \xxxSubParagraphStar
      \xxxSubParagraphNoStar
  }
  \newcommand{\xxxSubParagraphStar}[1]{\oldsubparagraph*{#1}\mbox{}}
  \newcommand{\xxxSubParagraphNoStar}[1]{\oldsubparagraph{#1}\mbox{}}
\fi
\makeatother


\usepackage{longtable,booktabs,array}
\newcounter{none} % for unnumbered tables
\usepackage{calc} % for calculating minipage widths
% Correct order of tables after \paragraph or \subparagraph
\usepackage{etoolbox}
\makeatletter
\patchcmd\longtable{\par}{\if@noskipsec\mbox{}\fi\par}{}{}
\makeatother
% Allow footnotes in longtable head/foot
\IfFileExists{footnotehyper.sty}{\usepackage{footnotehyper}}{\usepackage{footnote}}
\makesavenoteenv{longtable}
\usepackage{graphicx}
\makeatletter
\newsavebox\pandoc@box
\newcommand*\pandocbounded[1]{% scales image to fit in text height/width
  \sbox\pandoc@box{#1}%
  \Gscale@div\@tempa{\textheight}{\dimexpr\ht\pandoc@box+\dp\pandoc@box\relax}%
  \Gscale@div\@tempb{\linewidth}{\wd\pandoc@box}%
  \ifdim\@tempb\p@<\@tempa\p@\let\@tempa\@tempb\fi% select the smaller of both
  \ifdim\@tempa\p@<\p@\scalebox{\@tempa}{\usebox\pandoc@box}%
  \else\usebox{\pandoc@box}%
  \fi%
}
% Set default figure placement to htbp
\def\fps@figure{htbp}
\makeatother



\ifLuaTeX
\usepackage[bidi=basic,shorthands=off]{babel}
\else
\usepackage[bidi=default,shorthands=off]{babel}
\fi
\ifLuaTeX
  \usepackage{selnolig} % disable illegal ligatures
\fi


\setlength{\emergencystretch}{3em} % prevent overfull lines

\providecommand{\tightlist}{%
  \setlength{\itemsep}{0pt}\setlength{\parskip}{0pt}}



 


\KOMAoption{captions}{tableheading}
\usepackage[T1]{fontenc}
\usepackage{helvet}
\renewcommand{\familydefault}{\sfdefault}
% Fusszeile: Titel links, Seitenzahl rechts (auch auf der Titelseite)
\usepackage{scrlayer-scrpage}
\clearpairofpagestyles
\ifoot{Arbeitsblatt 2 (Massenbilanz) -- Principles of Data Visualization}
\ofoot{\pagemark}
\pagestyle{scrheadings}
% Standard-Titelblock von Quarto unterdruecken (Titelseite wird unten selbst gebaut)
\AtBeginDocument{\let\maketitle\relax}
\makeatletter
\@ifpackageloaded{tcolorbox}{}{\usepackage[skins,breakable]{tcolorbox}}
\@ifpackageloaded{fontawesome5}{}{\usepackage{fontawesome5}}
\definecolor{quarto-callout-color}{HTML}{909090}
\definecolor{quarto-callout-note-color}{HTML}{0758E5}
\definecolor{quarto-callout-important-color}{HTML}{CC1914}
\definecolor{quarto-callout-warning-color}{HTML}{EB9113}
\definecolor{quarto-callout-tip-color}{HTML}{00A047}
\definecolor{quarto-callout-caution-color}{HTML}{FC5300}
\definecolor{quarto-callout-color-frame}{HTML}{acacac}
\definecolor{quarto-callout-note-color-frame}{HTML}{4582ec}
\definecolor{quarto-callout-important-color-frame}{HTML}{d9534f}
\definecolor{quarto-callout-warning-color-frame}{HTML}{f0ad4e}
\definecolor{quarto-callout-tip-color-frame}{HTML}{02b875}
\definecolor{quarto-callout-caution-color-frame}{HTML}{fd7e14}
\makeatother
\makeatletter
\@ifpackageloaded{caption}{}{\usepackage{caption}}
\AtBeginDocument{%
\ifdefined\contentsname
  \renewcommand*\contentsname{Inhaltsverzeichnis}
\else
  \newcommand\contentsname{Inhaltsverzeichnis}
\fi
\ifdefined\listfigurename
  \renewcommand*\listfigurename{Abbildungsverzeichnis}
\else
  \newcommand\listfigurename{Abbildungsverzeichnis}
\fi
\ifdefined\listtablename
  \renewcommand*\listtablename{Tabellenverzeichnis}
\else
  \newcommand\listtablename{Tabellenverzeichnis}
\fi
\ifdefined\figurename
  \renewcommand*\figurename{Abbildung}
\else
  \newcommand\figurename{Abbildung}
\fi
\ifdefined\tablename
  \renewcommand*\tablename{Tabelle}
\else
  \newcommand\tablename{Tabelle}
\fi
}
\@ifpackageloaded{float}{}{\usepackage{float}}
\floatstyle{ruled}
\@ifundefined{c@chapter}{\newfloat{codelisting}{h}{lop}}{\newfloat{codelisting}{h}{lop}[chapter]}
\floatname{codelisting}{Listing}
\newcommand*\listoflistings{\listof{codelisting}{Listingverzeichnis}}
\makeatother
\makeatletter
\makeatother
\makeatletter
\@ifpackageloaded{caption}{}{\usepackage{caption}}
\@ifpackageloaded{subcaption}{}{\usepackage{subcaption}}
\makeatother
\usepackage{bookmark}
\IfFileExists{xurl.sty}{\usepackage{xurl}}{} % add URL line breaks if available
\urlstyle{same}
% fallback for those not using the hyperref driver hyperxmp:
\makeatletter
\@ifundefined{xmpquote}{\newcommand{\xmpquote}[1]{#1}}{}
\makeatother
\hypersetup{
  pdftitle={Jahresbilanz der Schweizer Gletscher},
  pdfauthor={Elio Schmid},
  pdflang={de-CH},
  colorlinks=true,
  linkcolor={blue},
  filecolor={Maroon},
  citecolor={Blue},
  urlcolor={Blue},
  pdfcreator={LaTeX via pandoc}}


\title{Jahresbilanz der Schweizer Gletscher}
\usepackage{etoolbox}
\makeatletter
\providecommand{\subtitle}[1]{% add subtitle to \maketitle
  \apptocmd{\@title}{\par {\large #1 \par}}{}{}
}
\makeatother
\subtitle{Arbeitsblatt 2, zweiter Datensatz -- Principles of Data
Visualization, Herbstsemester 2026}
\author{Elio Schmid}
\date{19.10.2026}
\begin{document}
\maketitle


\setstretch{1.3}
\thispagestyle{scrheadings}
\vspace*{0.5cm}
\begin{center}
{\bfseries\fontsize{24}{30}\selectfont Jahresbilanz der Schweizer Gletscher\par}
\vspace{1cm}
{\bfseries\fontsize{16}{20}\selectfont Arbeitsblatt 2, zweiter Datensatz\par}
\vspace{0.3cm}
{\bfseries\fontsize{16}{20}\selectfont Principles of Data Visualization -- Herbstsemester 2026\par}
\vspace{0.8cm}
\end{center}

\begin{center}
\includegraphics[width=0.65\linewidth,height=\textheight,keepaspectratio]{Schmid_Elio_Arbeitsblatt_2_Massenbilanz_files/figure-pdf/titelbild-1.pdf}
\end{center}

\vfill
\noindent
\begin{tabular}{@{}p{4cm}p{10cm}@{}}
Autor: & Elio Schmid,\\
       & FH Graubünden, MSc Data Visualization\\[1em]
Dozentin: & PD Dr. rer. nat. Helena Jambor\\[1em]
Ort, Datum: & Bern, 19.10.2026
\end{tabular}
\newpage

\section{Erster Datensatz:
Volumenänderung}\label{erster-datensatz-volumenuxe4nderung}

Anfang Oktober habe ich einen Artikel von SRF über den Gletscherschwund
in der Schweiz gelesen. Der Artikel ist hier
\href{https://www.srf.ch/wissen/klimawandel/eisverlust-2026-massive-gletscherschmelze-so-viel-eis-verschwand-im-hitzesommer}{verlinkt.}
Darin wird auf Daten von Glamos - Glacier Monitoring in Switzerland
verwiesen. Das Thema interessiert mich insofern, als dass ich relativ
viel in den Bergen unterwegs bin. Auf der Website von Glamos finden wir
Daten zu Bestand, Längenänderung, Massenbilanz, Volumenänderung sowie
Szenarien. Die Daten unterscheiden sich in Stichtagsmessungen und
Zeitreihen. Es werden Flächen-, Ausdehnungs- und Geodaten erfasst. In
meiner Analyse möchte ich mich auf die Volumenänderung konzentrieren.
Ich kann mir aber vorstellen auch Bestandesdaten einfliessen zu lassen
oder Längenänderungen. Dies werde ich im Zuge der Auswertung wohl
feststellen. Der Datensatz, den ich analysieren möchte, lässt sich
\href{https://doi.glamos.ch/data/volumechange/volumechange_2025_r2025.html}{hier}
einsehen.

Nach der ersten Auseinandersetzung mit dem Datensatz, muss ich
festhalten, dass es doch schwerer ist diesen auszuwerten als vorerst
angenommen. Die Beobachtungszeiträume überschneiden sich und Aussagen zu
Entwicklungen über eine Zeitdistanz sind anspruchsvoll. Sie können
höchstens für einen gewissen Teil der Gletscher erfolgen.Aus diesem
Grund starte ich einen zweiten Versuch mit den Daten der
Gletschermassenbilanz (hydrologisches Jahr). Ich erhoffe mir davon
bessere Möglichkeiten für Auswertungen über die Zeit, weil die Messungen
jährlich erfolgen und sich nicht überlappen.Auch der Upload in
Datawrapper hat mir gleich gezeigt, dass Auswertungen über die Zeit mit
diesem Datensatz ``einfacher'' zu realisieren sind.

\section{Zweiter Datensatz: Massenbilanz pro
Hydrojahr}\label{zweiter-datensatz-massenbilanz-pro-hydrojahr}

\subsection{Datensatz}\label{datensatz}

Auch dieser Datensatz ist bei GLAMOS verfügbar.Eine Zeile beschreibt
\textbf{einen Gletscher in einem Hydrojahr} (jeweils von Oktober bis
September). Die Winterbilanz (Bw) ist der Zugewinn durch Schnee bis Ende
April, die Sommerbilanz (Bs) der Verlust danach, die Jahresbilanz (Ba)
die Summe der beiden. Eine Bilanz von −1000 mm w.e. bedeutet, dass der
Gletscher im Mittel über seine Fläche eine Wasserschicht von 1 m
verloren hat. Zusätzlich gibt es die Gleichgewichtslinie (ELA), also die
Höhe, oberhalb der im Jahr netto Masse zugewinnt, und den Anteil der
Fläche oberhalb dieser Linie (AAR).

{\def\LTcaptype{none} % do not increment counter
\begin{longtable}[]{@{}
  >{\raggedright\arraybackslash}p{(\linewidth - 4\tabcolsep) * \real{0.2051}}
  >{\raggedright\arraybackslash}p{(\linewidth - 4\tabcolsep) * \real{0.6538}}
  >{\raggedright\arraybackslash}p{(\linewidth - 4\tabcolsep) * \real{0.1410}}@{}}
\toprule\noalign{}
\begin{minipage}[b]{\linewidth}\raggedright
Variable
\end{minipage} & \begin{minipage}[b]{\linewidth}\raggedright
Bedeutung
\end{minipage} & \begin{minipage}[b]{\linewidth}\raggedright
Einheit
\end{minipage} \\
\midrule\noalign{}
\endhead
\bottomrule\noalign{}
\endlastfoot
glacier\_name & Name des Gletschers & -- \\
glacier\_id & ID gemäss Swiss Glacier Inventory & -- \\
date\_start & Beginn des Hydrojahres & yyyy-mm-dd \\
date\_end\_winter & Ende der Winterbeobachtung & yyyy-mm-dd \\
date\_end & Ende des Hydrojahres & yyyy-mm-dd \\
Bw & Winterbilanz (Gletschermittel) & mm w.e. \\
Bs & Sommerbilanz (Gletschermittel) & mm w.e. \\
Ba & Jahresbilanz (Bw + Bs) & mm w.e. \\
ELA & Höhe der Gleichgewichtslinie & m ü. M. \\
AAR & Anteil der Fläche oberhalb der Gleichgewichtslinie & \% \\
area\_km2 & Gletscherfläche im Jahr & km² \\
h\_min & Tiefster Punkt des Gletschers & m ü. M. \\
h\_max & Höchster Punkt des Gletschers & m ü. M. \\
observer & Datenquelle bzw. Referenz & -- \\
\end{longtable}
}

\subsection{Daten einlesen}\label{daten-einlesen}

Die Datei hat 9 Kopfzeilen (Lizenz, Titel, ausführliche Namen,
Kurznamen, Einheiten), die Daten beginnen in Zeile 10. Die Spalte
\texttt{observer} enthält teils Kommas (z. B. «GLAMOS (2025), Huss 2015
(\ldots)»). Deshalb lese ich jede Zeile als Text ein und trenne nur nach
den ersten 13 Kommas, den Rest behalte ich als Quelle.

\begin{verbatim}
Rows: 1,239
Columns: 14
$ glacier_name    <chr> "Silvrettagletscher", "Silvrettagletscher", "Silvretta~
$ glacier_id      <chr> "A10g-05", "A10g-05", "A10g-05", "A10g-05", "A10g-05",~
$ date_start      <chr> "1914-10-01", "1915-10-01", "1916-10-01", "1917-10-01"~
$ date_end_winter <chr> "1915-04-30", "1916-04-30", "1917-04-30", "1918-04-30"~
$ date_end        <chr> "1915-09-30", "1916-09-30", "1917-09-30", "1918-09-30"~
$ Bw              <chr> "1503", "1415", "1872", "1075", "1862", "2467", "815",~
$ Bs              <chr> "-1918", "-1034", "-1900", "-1489", "-840", "-1884", "~
$ Ba              <chr> "-415", "381", "-28", "-414", "1022", "583", "-2546", ~
$ ELA             <chr> "2795", "2625", "2725", "2805", "2445", "2635", "3151"~
$ AAR             <chr> "45", "75", "60", "42", "98", "77", "0", "78", "80", "~
$ area_km2        <chr> "4.06687", "3.55563", "3.55563", "3.55563", "3.55563",~
$ h_min           <chr> "2406", "2414", "2414", "2414", "2414", "2414", "2414"~
$ h_max           <chr> "3185", "3091", "3091", "3091", "3091", "3091", "3091"~
$ observer        <chr> "GLAMOS (2025), Huss 2015 (https://doi.org/10.3189/201~
\end{verbatim}

\subsection{Bereinigung}\label{bereinigung}

Die Daten sind bereits sauber: keine fehlenden Werte, ISO-Datum, keine
Platzhalter. Es sind nur noch wenige Schritte nötig: Typen setzen, das
Jahr (Jahr des Periodenendes) und die Dauer berechnen und die Bilanzen
in m w.e. umrechnen, damit sie sich besser lesen lassen.

\subsection{Datenqualität}\label{datenqualituxe4t}

\begin{verbatim}
# A tibble: 0 x 2
# i 2 variables: variable <chr>, n_na <int>
\end{verbatim}

\begin{verbatim}
# A tibble: 0 x 3
# i 3 variables: glacier_id <chr>, jahr <dbl>, n <int>
\end{verbatim}

\begin{verbatim}
   Min. 1st Qu.  Median    Mean 3rd Qu.    Max. 
  364.0   364.0   364.0   364.3   365.0   365.0 
\end{verbatim}

\begin{verbatim}
[1] 0
\end{verbatim}

\begin{verbatim}
# A tibble: 1 x 3
  aar_null ela_ueber_gipfel extrapoliert
     <int>            <int>        <int>
1      117              153          163
\end{verbatim}

\begin{verbatim}
# A tibble: 6 x 2
  observer                                                                  n
  <chr>                                                                 <int>
1 GLAMOS (2025), Huss 2015 (https://doi.org/10.3189/2015JoG15J015)        730
2 VAW / ETH Zuerich, Huss et al. 2015, Journal.of Glaciology              226
3 GLAMOS (2025)                                                            88
4 GLAMOS (2025), Fischer 2016 (https://doi.org/10.5194/tc-10-1279-2016)    40
5 VAW / ETH Zuerich, O  Langenegger, unpublished                           35
6 GLAMOS (2025), Zekollari 2018 (https://doi.org/10.1017/jog.2018.18)      24
\end{verbatim}

Die Beobachtungen dazu:

\begin{itemize}
\tightlist
\item
  Es gibt keine fehlenden Werte und pro Gletscher und Jahr genau eine
  Zeile. Alle Hydrojahre dauern 364 oder 365 Tage. Die fünf auffälligen
  Gletscherjahre aus der Höhenstufen-Version (5 Tage, −14 Tage, 732
  Tage, 291 und 297 Tage) kommen hier nicht vor. Der einzige Unterschied
  bei den Gletscherjahren ist eine zusätzliche Zeile (Ofentalgletscher
  1965).
\item
  ELA und AAR sind nicht immer gemessene Werte. Bei 117 Gletscherjahren
  ist die AAR 0 (kein Nährgebiet), bei insgesamt 163 liegt die ELA über
  dem höchsten Punkt des Gletschers oder die AAR ist 0. Diese ELA-Werte
  sind extrapoliert und werden in den Grafiken getrennt markiert. 2022
  betrifft das 22 von 28 Gletschern.
\item
  Die Spalte \texttt{observer} ist hier die Datenquelle. Bei 956 von
  1'239 Zeilen wird Huss et al.~2015 genannt (in zwei Schreibweisen),
  eine Studie, die für viele Gletscher die Reihen rekonstruiert hat. Die
  Werte sind also nicht alle direkte Messungen, und die Quelle ist ein
  Hinweis auf unterschiedliche Datenqualität.
\end{itemize}

\subsection{Basiskennzahlen}\label{basiskennzahlen}

\begin{verbatim}
# A tibble: 1 x 6
  anzahl_zeilen anzahl_gletscher jahr_von jahr_bis median_Ba_m anteil_negativ
          <int>            <int>    <dbl>    <dbl>       <dbl>          <dbl>
1          1239               43     1885     2025      -0.559          0.748
\end{verbatim}

\begin{verbatim}
# A tibble: 12 x 5
   glacier_name               von   bis n_jahre groesste_luecke
   <chr>                    <dbl> <dbl>   <int>           <dbl>
 1 Grosser Aletschgletscher  1915  2025     111               1
 2 Silvrettagletscher        1915  2025     111               1
 3 Claridenfirn              1915  2025     109               3
 4 Allalingletscher          1956  2025      70               1
 5 Hohlaubgletscher          1956  2025      70               1
 6 Schwarzberggletscher      1956  2025      70               1
 7 Griesgletscher            1962  2025      64               1
 8 Glacier du Giétro         1967  2025      59               1
 9 Rhonegletscher            1885  2025      44              98
10 Tällibodengletscher       1956  1996      41               1
11 Chessjengletscher         1956  1995      40               1
12 Limmerngletscher          1948  1985      38               1
\end{verbatim}

\begin{verbatim}
# A tibble: 8 x 2
   jahr     n
  <dbl> <int>
1  1920     3
2  1940     3
3  1960    13
4  1980    12
5  2000    10
6  2010    19
7  2020    30
8  2025    27
\end{verbatim}

\subsection{Erste Übersicht}\label{erste-uxfcbersicht}

\begin{figure}

\centering{

\pandocbounded{\includegraphics[keepaspectratio]{Schmid_Elio_Arbeitsblatt_2_Massenbilanz_files/figure-pdf/fig-hist-1.pdf}}

}

\caption{\label{fig-hist}Verteilung der Jahresbilanz pro Gletscher und
Jahr}

\end{figure}%

\begin{figure}

\centering{

\pandocbounded{\includegraphics[keepaspectratio]{Schmid_Elio_Arbeitsblatt_2_Massenbilanz_files/figure-pdf/fig-erste-zeit-1.pdf}}

}

\caption{\label{fig-erste-zeit}Erste Übersicht: Jahresbilanz aller
Gletscherjahre über die Zeit}

\end{figure}%

Die Punktwolke zeigt schon auf den ersten Blick einen Trend zu
negativeren Werten, aber auch, dass die Anzahl Punkte pro Jahr stark
schwankt. Ein einfacher Mittelwert pro Jahr würde deshalb Veränderungen
der Gletscherauswahl mit echten Veränderungen vermischen. Das greife ich
in Ansatz 1 auf.

\subsection{Welche Fragestellungen möchte ich mit diesem Datensatz
beantworten?}\label{welche-fragestellungen-muxf6chte-ich-mit-diesem-datensatz-beantworten}

Bereits bei meinem ersten Versuch wollte ich eine Entwicklung über die
Zeit darstellen.Das möchte ich nun auch mit diesem Datensatz versuchen.
Zudem können wir einzelne Gletscher aus unterschiedlichen Regionen der
Schweiz vergleichen. Wenn möglich, möchte ich Zusammenhänge von z.B. der
Höhe und der Entwicklung der Massenbilanz darstellen.

\subsection{Visualisierungsansätze}\label{visualisierungsansuxe4tze}

Nachfolgend erste Visualisierungsansätze. Der Datensatz beinhaltet
insgesamt 9 Gletscher, die eine fast lückenlose Messreihe ab 1990
aufweisen. Diese werden für eine Aussage über die Zeit verwendet.Bei den
nachfolgenden Visualisierungen handelt es sich um erste Würfe. Farbliche
Ausgestaltung und ergänzende Texte werden noch vertiefter bearbeitet

\begin{verbatim}
[1] "Silvrettagletscher"       "Claridenfirn"            
[3] "Grosser Aletschgletscher" "Griesgletscher"          
[5] "Schwarzberggletscher"     "Allalingletscher"        
[7] "Hohlaubgletscher"         "Glacier du Giétro"       
[9] "Ghiacciaio del Basòdino" 
\end{verbatim}

\subsubsection{Ansatz 1: Heatmap aller Gletscher über die
Zeit}\label{ansatz-1-heatmap-aller-gletscher-uxfcber-die-zeit}

\textbf{Frage:} Welche Gletscher wurden wann gemessen, und welche Jahre
waren für alle Gletscher gleichzeitig extrem?

\textbf{Aufbau der Visualisierung:} Muster über Gletscher und Zeit
erkennen, Datenlücken sichtbar machen. Rechtecke in einem Raster. Die
x-Achse zeigt das Jahr, die y-Achse den Gletscher (sortiert nach dem
ersten Messjahr), die Farbe die Jahresbilanz auf einer divergierenden
Skala mit Mitte bei null. Weisse Felder bedeuten, dass kein Wert
vorliegt. So sehen wir Datenlücken und Werte in derselben Grafik.

\begin{figure}

\centering{

\pandocbounded{\includegraphics[keepaspectratio]{Schmid_Elio_Arbeitsblatt_2_Massenbilanz_files/figure-pdf/fig-ansatz2-1.pdf}}

}

\caption{\label{fig-ansatz2}Ansatz 1: Jahresbilanz pro Gletscher und
Jahr (Farbe: Rot = Verlust, Blau = Zuwachs)}

\end{figure}%

\textbf{Erkenntnis:} Aus der Heatmap können mehrere Informationen
gleichzeitig ausgelesen werden. Erstens zeigt sie die Datenlage: Drei
Gletscher (Silvretta, Aletsch, Claridenfirn) wurden seit 1915 gemessen,
eine grössere Gruppe ab Mitte der 1950er-Jahre, die meisten anderen erst
seit etwa 2010. Der Rhonegletscher hat zwei getrennte Blöcke (1885 bis
1909 und ab 2007). Zweitens sieht man, dass einzelne Jahre bei fast
allen Gletschern gleichzeitig dunkelrot sind (2022, 2003), während die
Blautöne der 1960er- und 1970er-Jahre in den letzten Jahrzehnten fast
verschwunden sind.In dieser Version müssen die fehlenden Werte farblich
noch überarbeitet werden. Die Unterscheidung zwischen einer neutralen
Massenbilanz (0 Wert) und einem fehlenden Wert ist schwer zu
machen.Überprüft wird zudem ob die Skala passend ist.

\paragraph{Iteration 2: Heatmap ab 2000 und ab
1990}\label{iteration-2-heatmap-ab-2000-und-ab-1990}

\textbf{Anpassungen:} Die erste Version enthielt viele leere Felder,
weil einzelne Gletscher erst spät oder mit Lücken gemessen wurden. In
dieser Iteration zeige ich nur Gletscher mit einer fast lückenlosen
Reihe: für die Grafik ab 2000 alle Gletscher mit mindestens 20 von 26
Jahren, für die Grafik ab 1990 die neun Gletscher des konstanten
Samples. Fehlende Werte sind neu \textbf{grau} dargestellt, eine
ausgeglichene Bilanz (0) bleibt \textbf{weiss}. Zusätzlich sind die
Gletscher nach der mittleren Bilanz sortiert (stärkster Verlust oben),
damit Muster sichtbar werden.

\begin{figure}

\centering{

\pandocbounded{\includegraphics[keepaspectratio]{Schmid_Elio_Arbeitsblatt_2_Massenbilanz_files/figure-pdf/fig-iter2-heatmap-2000-1.pdf}}

}

\caption{\label{fig-iter2-heatmap-2000}Iteration 2: Jährliche
Massenbilanz ausgewählter Schweizer Gletscher seit 2000 (grau = kein
Messwert, weiss = ausgeglichene Bilanz)}

\end{figure}%

\begin{figure}

\centering{

\pandocbounded{\includegraphics[keepaspectratio]{Schmid_Elio_Arbeitsblatt_2_Massenbilanz_files/figure-pdf/fig-iter2-heatmap-1990-1.pdf}}

}

\caption{\label{fig-iter2-heatmap-1990}Iteration 2: Jährliche
Massenbilanz der neun Gletscher des konstanten Samples seit 1990 (grau =
kein Messwert, weiss = ausgeglichene Bilanz)}

\end{figure}%

\textbf{Erkenntnis:} Mit den ausgewählten Gletschern und der Sortierung
nach mittlerer Bilanz sind die Muster deutlich besser lesbar. Seit 2000
sind von 305 Gletscherjahren nur 27 (9 \%) positiv, und nur vier liegen
über +0.5 m w.e. (2001 und 2007). Die Jahre 2003 und 2022 sind bei allen
Gletschern dunkelrot, 2023 ist ebenfalls bei fast allen stark negativ.
In der Grafik ab 1990 sieht man mehr blaue Felder (42 von 320
Gletscherjahren, 13 \%), vor allem Anfang der 1990er-Jahre und um 2001.
Seit 2011 sind nur noch 4 \% der Gletscherjahre positiv, zwischen 1990
und 2010 waren es rund 20 \%. In 20 der 36 Jahre sind alle neun
Gletscher negativ. Die grauen Felder (fehlende Werte) sind nun klar von
den weissen (ausgeglichene Bilanz) unterscheidbar: Es sind sieben Felder
bei der Grafik ab 2000 und vier bei der Grafik ab 1990. Die Skala bei ±3
m w.e. reicht für die Grafik ab 1990 fast überall aus, nur einzelne
Werte im Jahr 2022 (bis −3.9) werden abgeschnitten. Auf der positiven
Seite wird die Skala nur bis +1.3 m w.e. genutzt, deshalb sind die
Blautöne blass.

\paragraph{Iteration 3: Heatmap aller Gletscher seit
1990}\label{iteration-3-heatmap-aller-gletscher-seit-1990}

\textbf{Anpassungen:} Nun alle Gletscher, die seit 1990 gemessen wurden,
nicht nur eine Auswahl. Das Farbkonzept (Rot = Verlust, Blau = Zuwachs,
Weiss = ausgeglichen, Grau = kein Messwert) bleibt gleich. Die Gletscher
sind neu nach dem ersten Messjahr sortiert (früheste Messreihe oben),
weil die Sortierung nach mittlerer Bilanz bei unterschiedlichen
Messzeiträumen nicht vergleichbar ist. Die Grafik ist höher, damit alle
Gletschernamen lesbar sind.

\begin{figure}

\centering{

\pandocbounded{\includegraphics[keepaspectratio]{Schmid_Elio_Arbeitsblatt_2_Massenbilanz_files/figure-pdf/fig-iter3-heatmap-alle-1990-1.pdf}}

}

\caption{\label{fig-iter3-heatmap-alle-1990}Iteration 3: Jährliche
Massenbilanz aller Schweizer Gletscher seit 1990, sortiert nach erstem
Messjahr (grau = kein Messwert, weiss = ausgeglichene Bilanz)}

\end{figure}%

\textbf{Erkenntnis:} Die Sortierung nach erstem Messjahr macht den
Ausbau des Messnetzes sichtbar: Seit 1990 wurden 38 Gletscher gemessen,
aber nur 10 davon schon in den frühen 1990er-Jahren. Die Zahl der
Gletscher mit Messwert steigt von 12 (2005) auf 19 (2010) und 29 (2015).
Gut die Hälfte der 1'368 möglichen Felder ist deshalb grau (701), vor
allem links unten. Die Grafik eignet sich darum weniger für den
Vergleich einzelner Gletscher als für das Gesamtbild: Von 667
Gletscherjahren sind nur 61 (9 \%) positiv, und nur zwei liegen über +1
m w.e. Bis 2010 waren es 17 \% positive Gletscherjahre, seit 2011 knapp
5 \%, seit 2015 weniger als 1 \%. In den Jahren 2003, 2022 und 2023 sind
alle gemessenen Gletscher negativ (2022: alle 28, bis −4.1 m w.e.). Die
senkrechten dunkelroten Streifen bei 2003 und 2022 zeigen, dass es sich
um Jahreseffekte handelt, die alle Gletscher gleichzeitig treffen,
unabhängig von der Länge der Messreihe. Die oberen neun Gletscher mit
der längsten Reihe zeigen zusätzlich den Wechsel von helleren Jahren
(1990er) zu fast durchgehend roten Jahren seit 2010.

\subsubsection{Ansatz 2: Jahresbilanz im Zeitverlauf mit konstantem
Sample}\label{ansatz-2-jahresbilanz-im-zeitverlauf-mit-konstantem-sample}

\textbf{Frage:} Wie hat sich die Jahresbilanz verschiedener Schweizer
Gletscher seit 1990 entwickelt?

\textbf{Aufbau der Visualisierung:} Für den Vergleich über die Zeit ist
ein Liniendiagramm sinnvoll. Die Zeit liegt auf der x-Achse, die
Jahresbilanz auf der y-Achse Die dünnen grauen Linien sind die einzelnen
Gletscher, die dicke schwarze Linie ist der Mittelwert über die 9 zuvor
genannten Gletscher, die in praktisch allen Jahren seit 1990 gemessen
wurden. Die Linien bilden dadurch die echte Veränderung ab und nicht
eine Veränderung die ensteht wegen der Veränderung der Anzahl gemessenen
Gletscher.

\begin{figure}

\centering{

\pandocbounded{\includegraphics[keepaspectratio]{Schmid_Elio_Arbeitsblatt_2_Massenbilanz_files/figure-pdf/fig-ansatz1-1.pdf}}

}

\caption{\label{fig-ansatz1}Ansatz 1: Jahresbilanz seit 1990, einzelne
Gletscher (grau) und Mittel des konstanten Samples (schwarz)}

\end{figure}%

\textbf{Erkenntnis:} In den 36 Jahren seit 1990 ist das Mittel nur in
drei Jahren (darunter 1993 und 2001) knapp positiv.Besonders auffällig
ist das Jahr 2022, welches ein sehr negatives Resultat aufweist.

\paragraph{Iteration 2: Zeitverlauf mit beschrifteten
Gletschern}\label{iteration-2-zeitverlauf-mit-beschrifteten-gletschern}

\textbf{Anpassungen:} Die Grafik zeigt nur noch die neun Gletscher des
konstanten Samples, damit die Linien genau dem Mittelwert entsprechen
und alle direkt beschriftet werden können (bei allen rund 30 Gletschern
seit 1990 wären die Beschriftungen am Linienende nicht mehr lesbar). Die
Nulllinie ist neu durchgezogen und dunkler, der Mittelwert ist farbig
und in der Legende benannt, die zwei auffälligen Jahre sind beschriftet.
Die Namen der Gletscher stehen am Ende der Linien in Grau und kursiv.

\begin{figure}

\centering{

\pandocbounded{\includegraphics[keepaspectratio]{Schmid_Elio_Arbeitsblatt_2_Massenbilanz_files/figure-pdf/fig-iter2-zeitverlauf-1.pdf}}

}

\caption{\label{fig-iter2-zeitverlauf}Iteration 2: Entwicklung der
jährlichen Massenbilanz ausgewählter Schweizer Gletscher seit 1990}

\end{figure}%

\textbf{Erkenntnis:} Der farbige Mittelwert hebt sich nun klar von den
grauen Gletscherlinien ab, und die durchgezogene Nulllinie macht
sichtbar, dass das Mittel nur in drei Jahren über null liegt (1993, 1995
und 2001). Die beiden Extremjahre sind beschriftet: 2003 (Hitzesommer,
Mittel −2.1 m w.e., alle neun Gletscher zwischen −2.7 und −1.2) und 2022
(Rekordjahr, Mittel −3.1 m w.e., zwischen −3.9 und −2.4). Seit 2023 hat
sich die Bilanz etwas erholt (−2.4, −1.1, −1.5 m w.e. in den Jahren 2023
bis 2025), bleibt aber klar negativ. Alle neun Gletscher lassen sich am
Linienende direkt beschriften. Bei allen rund 30 Gletschern seit 1990
wäre das nicht mehr lesbar. Die Gletscher verlaufen weitgehend parallel,
der Unterschied zwischen ihnen ist im Vergleich zu den Schwankungen
zwischen den Jahren klein. Am stärksten ist der Verlust im Mittel beim
Griesgletscher (−1.4 m w.e. pro Jahr seit 1990), am schwächsten beim
Allalingletscher (−0.6).

\subsubsection{Ansatz 3: Winter- und Sommerbilanz
getrennt}\label{ansatz-3-winter--und-sommerbilanz-getrennt}

\textbf{Frage:} Entsteht der Massenverlust durch weniger Schnee im
Winter oder durch mehr Schmelze im Sommer?

\textbf{Aufbau der Visualisierung} Balken vom Nullpunkt aus. Die
Winterbilanz (Akkumulation) zeigt nach oben, die Sommerbilanz
(Massenverlust) nach unten, jeweils auf derselben y-Skala. Die schwarze
Linie ist die Jahresbilanz, die Summe der beiden. Alle drei Werte
stammen vom konstanten Sample aus Ansatz 2. Die Werte sind so direkt
Verlgeichbar.

\begin{figure}

\centering{

\pandocbounded{\includegraphics[keepaspectratio]{Schmid_Elio_Arbeitsblatt_2_Massenbilanz_files/figure-pdf/fig-ansatz3-1.pdf}}

}

\caption{\label{fig-ansatz3}Ansatz 3: Winterbilanz (oben), Sommerbilanz
(unten) und Jahresbilanz (Linie), Mittel des konstanten Samples}

\end{figure}%

\textbf{Erkenntnis:} Die Aufteilung zeigt, dass die Winterbilanz im
Mittel stabil blieb, während die Sommerbilanz negativer wird. Die
Jahresbilanz folgt fast ausschliesslich der Sommerbilanz (Korrelation
0.92, mit der Winterbilanz nur 0.35). Auch das Jahr 2022 erklärt sich
so: Die Schneemenge im Winter war mit 0.86 m w.e. unterdurchschnittlich,
vor allem aber war die Sommerschmelze mit −4.0 m w.e. die höchste im
Sample.

\paragraph{Iteration 2: Winter- und Sommerbilanz mit beschriftetem
Mittelwert}\label{iteration-2-winter--und-sommerbilanz-mit-beschriftetem-mittelwert}

\textbf{Anpassungen:} Die Jahresbilanz (Mittel der neun Gletscher) ist
neu farbig, in der Legende benannt und zusätzlich direkt am Linienende
beschriftet. Ein Titel nennt Inhalt und Zeitraum.

\begin{figure}

\centering{

\pandocbounded{\includegraphics[keepaspectratio]{Schmid_Elio_Arbeitsblatt_2_Massenbilanz_files/figure-pdf/fig-iter2-saisons-1.pdf}}

}

\caption{\label{fig-iter2-saisons}Iteration 2: Winter-, Sommer- und
Jahresbilanz (Mittel der neun Gletscher) seit 1990}

\end{figure}%

\textbf{Erkenntnis:} Der Mittelwert ist nun farbig, im Diagramm
beschriftet und in der Legende erklärt. An der Aussage ändert sich
nichts: Die Winterbilanz bleibt über den ganzen Zeitraum ungefähr gleich
hoch, während die Sommerbilanz tiefer wird. Die zwei Extremjahre sind
direkt in den Balken sichtbar: 2003 (Sommerbilanz −3.5 m w.e.) und 2022
(−4.0 m w.e., bei gleichzeitig unterdurchschnittlicher Winterbilanz von
0.9 m w.e.).

\subsection{Zwischenfazit}\label{zwischenfazit}

\textbf{Einschätzung des Datensatzes:} Für Zeitauswertungen eignet sich
dieser Datensatz deutlich besser als die Volumenänderung. Die Perioden
sind jährlich, ein Gletscherjahr hat genau einen Wert. Der Preis ist
eine kleine und zeitlich ungleiche Auswahl: Nur 43 Gletscher sind
enthalten, nur drei davon seit 1915, und die Zahl der Gletscher pro Jahr
schwankt zwischen 3 und 31. Eine Aussage für «die Schweizer Gletscher»
ist damit nicht möglich.

\textbf{Offene Punkte:}

\begin{itemize}
\tightlist
\item
  Ein kumulativer Verlauf (Summe der Jahresbilanzen ab 1990) könnte die
  Entwicklung noch anschaulicher machen.
\item
  Ansatz 5 könnte für weitere Gletscher oder als Animation über die
  Jahre getestet werden.
\item
  Die Farbwahl der Heatmap (Rot--Blau) sollte ich noch auf
  Farbsehschwäche prüfen.
\item
  Die beiden Datensätze (Volumenänderung und Massenbilanz) liessen sich
  verknüpfen, um die geodätischen mit den glaziologischen Werten zu
  vergleichen.
\item
  Einzelne Gletscher mit dem absoluten Massenverlust zeigen. Allenfalls
  die drei Gletscher für welche Messdaten seit 1915 bestehen.
\item
  Karte mit einzelnen Angaben zu Gletschern darstellen.
\item
  Grafik wie bei Gun-Death-Florida für die Gesamtbilanz über die Jahre.
\item
  Bei Grafik 1 Gletscher der Stichprobe markieren.
\item
  Small Multiples verwenden
\end{itemize}

\subsection{Peer-Feedback und nächste
Iteration}\label{peer-feedback-und-nuxe4chste-iteration}

\begin{tcolorbox}[enhanced jigsaw, arc=.35mm, bottomrule=.15mm, bottomtitle=1mm, breakable, colback=white, colbacktitle=quarto-callout-note-color!10!white, colframe=quarto-callout-note-color-frame, coltitle=black, left=2mm, leftrule=.75mm, opacityback=0, opacitybacktitle=0.6, rightrule=.15mm, title=\textcolor{quarto-callout-note-color}{\faInfo}\hspace{0.5em}{Hinweis}, titlerule=0mm, toprule=.15mm, toptitle=1mm]

Wird nach der Vorlesung mit dem Feedback der anderen Studierenden
ergänzt.

\end{tcolorbox}




\end{document}
